\documentclass{article}
\usepackage{amsmath,amssymb,amsthm}
\usepackage{algorithm,algorithmic}
\usepackage{bm}
\usepackage{hyperref}
\usepackage{enumitem}
\usepackage{geometry}
\geometry{margin=1in}
\newtheorem{theorem}{Theorem}
\newtheorem{lemma}{Lemma}
\newcommand{\E}{\mathbb{E}}
\newcommand{\R}{\mathbb{R}}
\newcommand{\norm}[1]{\left\| #1 \right\|}
\newcommand{\ip}[2]{\left\langle #1,#2 \right\rangle}
\newcommand{\Q}{\mathcal{Q}}
\begin{document}

\section*{Quantized Stochastic Gradient Descent (QSGD)}

\section*{Mathematical Formulation}
Let $f:\R^d\to\R$ be a convex and $L$-smooth objective.  
At iteration $t$ we draw an i.i.d.\ data sample $\xi_t$ and compute a stochastic gradient 
\[
g_t \;:=\; \nabla f(w_t;\xi_t).
\]
Instead of sending $g_t$ we apply a (possibly stochastic) \emph{quantization operator}
\[
\Q(g_t)\in\R^d,\qquad\text{such that}\qquad \E[\Q(g_t)\mid g_t]=g_t .
\]
The master (or the next worker) updates the iterate by
\begin{equation}\label{eq:update}
w_{t+1}=w_t-\eta_t\,\Q(g_t),
\end{equation}
where $\{\eta_t\}_{t\ge 0}$ is a stepsize (learning‑rate) schedule.

\section*{Pseudocode}
\begin{algorithm}[H]
\caption{Quantized Stochastic Gradient Descent (QSGD)}
\begin{algorithmic}[1]
\STATE \textbf{Input:} initial point $w_0\in\R^d$, stepsize schedule $\{\eta_t\}$, quantizer $\Q(\cdot)$
\FOR{$t=0,1,\dots,T-1$}
    \STATE Sample $\xi_t\sim\mathcal{D}$  \hfill\COMMENT{draw a random data point}
    \STATE Compute stochastic gradient $g_t = \nabla f(w_t;\xi_t)$
    \STATE Quantize: $\tilde g_t = \Q(g_t)$
    \STATE Update: $w_{t+1}= w_t-\eta_t \tilde g_t$ \hfill\COMMENT{Eq.~\eqref{eq:update}}
\ENDFOR
\STATE \textbf{Output:} $\bar w_T = \frac1T\sum_{t=0}^{T-1} w_t$ (optional averaging)
\end{algorithmic}
\end{algorithm}

\section*{Dry Run Example}
Consider the scalar quadratic
\[
f(w)= (w-3)^2, \qquad \nabla f(w)=2(w-3).
\]
Take $w_0=0$, constant stepsize $\eta=0.1$, and use the identity quantizer $\Q(g)=g$ (so $\tilde g_t=g_t$).  

\begin{center}
\begin{tabular}{c|c|c|c}
$t$ & $w_t$ & $g_t=2(w_t-3)$ & $w_{t+1}=w_t-\eta g_t$ \\ \hline
0 & $0$ & $-6$ & $0-0.1(-6)=0.6$ \\
1 & $0.6$ & $2(0.6-3)=-4.8$ & $0.6-0.1(-4.8)=1.08$ \\
2 & $1.08$ & $2(1.08-3)=-3.84$ & $1.08-0.1(-3.84)=1.464$ \\
\end{tabular}
\end{center}
Corresponding objective values:
\[
f(0)=9,\; f(0.6)=5.76,\; f(1.08)=3.6864,\; f(1.464)=2.3415.
\]
The iterates move toward the minimiser $w^\star=3$.

\newpage

\section*{Assumptions}
\begin{enumerate}[label=(A\arabic*)]
\item \textbf{Convexity and smoothness.} $f(\cdot)$ is convex and $L$-smooth:
\[
f(y)\le f(x)+\ip{\nabla f(x)}{y-x}+\frac{L}{2}\norm{y-x}^2,\qquad\forall x,y\in\R^d.
\]
\item \textbf{Unbiased stochastic gradient.} For all $t$,
\[
\E[g_t\mid w_t]=\nabla f(w_t),\qquad \E[\norm{g_t}^2\mid w_t]\le G^2.
\]
\item \textbf{Unbiased quantization with bounded variance.} There exists a constant $q\ge 0$ such that
\[
\E[\Q(g_t)\mid g_t]=g_t,\qquad 
\E\!\left[\norm{\Q(g_t)-g_t}^2\mid g_t\right]\le q\,\norm{g_t}^2 .
\]
\item \textbf{Stepsize.} The stepsize sequence is non‑increasing and satisfies
\[
\sum_{t=0}^{\infty}\eta_t=\infty,\qquad
\sum_{t=0}^{\infty}\eta_t^{2}<\infty .
\]
A typical choice is $\eta_t=\frac{c}{\sqrt{t+1}}$ with $c>0$.
\end{enumerate}

\section*{Main Convergence Theorem}
\begin{theorem}\label{thm:conv}
Under (A1)–(A4) and for the averaged iterate $\bar w_T=\frac1T\sum_{t=0}^{T-1}w_t$, the Quantized SGD iterates satisfy
\[
\E\!\bigl[f(\bar w_T)-f(w^\star)\bigr]\;\le\;
\frac{\norm{w_0-w^\star}^2}{2\sum_{t=0}^{T-1}\eta_t}
\;+\;
\frac{(1+q)G^{2}}{2}\,\frac{\sum_{t=0}^{T-1}\eta_t}{\sum_{t=0}^{T-1}\eta_t}.
\]
In particular, with the stepsize $\eta_t=\frac{c}{\sqrt{t+1}}$,
\[
\E\!\bigl[f(\bar w_T)-f(w^\star)\bigr]\;=\;O\!\left(\frac{1}{\sqrt{T}}\right).
\]
\end{theorem}

\section*{Theoretical Properties}
\begin{itemize}
\item \textbf{Stability.} The additional variance term $q\,\E[\norm{g_t}^2]$ appears only multiplicatively; if $q=0$ the bound reduces to the classic SGD bound.
\item \textbf{Hyper‑parameter sensitivity.} The constant $c$ in $\eta_t=c/\sqrt{t+1}$ must be chosen such that $c\le \frac{1}{L\sqrt{1+q}}$ to guarantee the descent inequality used in the proof.
\item \textbf{Comparison to baselines.} Compared with vanilla SGD (where $q=0$), quantized SGD suffers at most a factor $1+q$ slowdown in the asymptotic rate, while drastically reducing communication bandwidth.
\end{itemize}

\section*{Preliminaries (Definitions and Lemmas)}
\begin{lemma}[Three‑point identity for smooth functions]\label{lem:smooth}
For any $x,y\in\R^d$,
\[
f(y)\le f(x)+\ip{\nabla f(x)}{y-x}+\frac{L}{2}\norm{y-x}^2 .
\]
\end{lemma}
\begin{lemma}[Unbiasedness of the quantizer]\label{lem:unbiased}
For any random vector $g$,
\[
\E[\Q(g)\mid g]=g .
\]
\end{lemma}
\begin{lemma}[Variance decomposition]\label{lem:var}
For any random vector $Z$,
\[
\E[\norm{Z}^2]=\norm{\E[Z]}^2+\E\bigl[\norm{Z-\E[Z]}^2\bigr].
\]
\end{lemma}

\section*{Main Proof (Step‑by‑Step Derivation)}
We prove Theorem~\ref{thm:conv}.  
All expectations are taken with respect to the randomness of the data sample $\xi_t$ and the quantizer $\Q(\cdot)$ unless otherwise noted.

\subsection*{Step 1: One‑step progress}
From the update \eqref{eq:update} and expanding the squared distance to the optimum $w^\star$,
\begin{align}
\norm{w_{t+1}-w^\star}^2
&= \norm{w_t-\eta_t\Q(g_t)-w^\star}^2 \nonumber\\
&= \norm{w_t-w^\star}^2 -2\eta_t\ip{\Q(g_t)}{w_t-w^\star}
   +\eta_t^{2}\norm{\Q(g_t)}^{2}. \tag{Step 1}
\end{align}

\subsection*{Step 2: Take conditional expectation}
Condition on $\mathcal{F}_t:=\sigma(w_0,\dots,w_t)$.
Using Lemma~\ref{lem:unbiased} and the tower property,
\begin{align}
\E\bigl[ \norm{w_{t+1}-w^\star}^2 \mid \mathcal{F}_t \bigr]
&= \norm{w_t-w^\star}^2 
   -2\eta_t\ip{\E[\Q(g_t)\mid\mathcal{F}_t]}{w_t-w^\star}
   +\eta_t^{2}\E\!\bigl[\norm{\Q(g_t)}^{2}\mid\mathcal{F}_t \bigr] \nonumber\\
&= \norm{w_t-w^\star}^2 
   -2\eta_t\ip{ \nabla f(w_t)}{w_t-w^\star}
   +\eta_t^{2}\E\!\bigl[\norm{\Q(g_t)}^{2}\mid\mathcal{F}_t \bigr]. \tag{Step 2}
\end{align}

\subsection*{Step 3: Upper‑bound the quantized variance}
Apply Lemma~\ref{lem:var} with $Z=\Q(g_t)$ and Lemma~\ref{lem:unbiased}:
\[
\E\!\bigl[\norm{\Q(g_t)}^{2}\mid\mathcal{F}_t \bigr]
 = \norm{\E[\Q(g_t)\mid\mathcal{F}_t]}^{2}
   +\E\!\bigl[\norm{\Q(g_t)-\E[\Q(g_t)\mid\mathcal{F}_t]}^{2}\mid\mathcal{F}_t \bigr]
 = \norm{g_t}^{2}+ \E\!\bigl[\norm{\Q(g_t)-g_t}^{2}\mid\mathcal{F}_t \bigr].
\]
By assumption (A3),
\[
\E\!\bigl[\norm{\Q(g_t)-g_t}^{2}\mid\mathcal{F}_t \bigr]\le q\,\norm{g_t}^{2}.
\]
Hence
\[
\E\!\bigl[\norm{\Q(g_t)}^{2}\mid\mathcal{F}_t \bigr]\le (1+q)\norm{g_t}^{2}. \tag{Step 3}
\]

\subsection*{Step 4: Substitute Step 3 into Step 2}
\begin{align}
\E\bigl[ \norm{w_{t+1}-w^\star}^2 \mid \mathcal{F}_t \bigr]
&\le \norm{w_t-w^\star}^2 
   -2\eta_t\ip{ \nabla f(w_t)}{w_t-w^\star}
   +\eta_t^{2}(1+q)\norm{g_t}^{2}. \tag{Step 4}
\end{align}

\subsection*{Step 5: Relate inner product to function values (convexity)}
Convexity of $f$ yields
\[
f(w_t)-f(w^\star)\le \ip{\nabla f(w_t)}{w_t-w^\star}. \tag{Step 5}
\]

\subsection*{Step 6: Upper‑bound $\norm{g_t}^2$}
Using (A2) and the law of total expectation,
\[
\E\!\bigl[\norm{g_t}^{2}\mid\mathcal{F}_t\bigr]
= \E\!\bigl[\norm{g_t-\nabla f(w_t)}^{2}\mid\mathcal{F}_t \bigr]
  +\norm{\nabla f(w_t)}^{2}
\le \sigma^{2}+ \norm{\nabla f(w_t)}^{2},
\]
where $\sigma^{2}:=G^{2}-\norm{\nabla f(w_t)}^{2}\le G^{2}$ is a uniform bound.
For simplicity we keep the coarse bound $\norm{g_t}^{2}\le G^{2}$ almost surely, which follows from (A2). Thus
\[
\E\!\bigl[\norm{g_t}^{2}\mid\mathcal{F}_t\bigr]\le G^{2}. \tag{Step 6}
\]

\subsection*{Step 7: Take full expectation}
Taking expectation of both sides of \eqref{Step 4} and using (Step 5)–(Step 6):
\begin{align}
\E\!\bigl[\norm{w_{t+1}-w^\star}^2\bigr]
&\le \E\!\bigl[\norm{w_t-w^\star}^2\bigr]
   -2\eta_t \E\!\bigl[f(w_t)-f(w^\star)\bigr]
   +\eta_t^{2}(1+q)G^{2}. \tag{Step 7}
\end{align}

\subsection*{Step 8: Rearrangement}
Rearrange (Step 7) to isolate the suboptimality term:
\begin{align}
2\eta_t \E\!\bigl[f(w_t)-f(w^\star)\bigr]
&\le \E\!\bigl[\norm{w_t-w^\star}^2\bigr]
   -\E\!\bigl[\norm{w_{t+1}-w^\star}^2\bigr]
   +\eta_t^{2}(1+q)G^{2}. \tag{Step 8}
\end{align}

\subsection*{Step 9: Sum over $t=0,\dots,T-1$}
Summing (Step 8) yields
\begin{align}
2\sum_{t=0}^{T-1}\eta_t \E\!\bigl[f(w_t)-f(w^\star)\bigr]
&\le \norm{w_0-w^\star}^2
   + (1+q)G^{2}\sum_{t=0}^{T-1}\eta_t^{2}. \tag{Step 9}
\end{align}
(The telescoping cancelation of $\E\!\bigl[\norm{w_{t+1}-w^\star}^2\bigr]$ occurs.)

\subsection*{Step 10: Use convexity of the average}
By convexity of $f$,
\[
f(\bar w_T)-f(w^\star)\le\frac{1}{\sum_{t=0}^{T-1}\eta_t}
      \sum_{t=0}^{T-1}\eta_t\bigl[f(w_t)-f(w^\star)\bigr].
\]
Taking expectation and applying (Step 9):
\begin{align}
\E\!\bigl[f(\bar w_T)-f(w^\star)\bigr]
&\le \frac{1}{2\sum_{t=0}^{T-1}\eta_t}
    \Bigl[\norm{w_0-w^\star}^2
    +(1+q)G^{2}\sum_{t=0}^{T-1}\eta_t^{2}\Bigr]. \tag{Step 10}
\end{align}

\subsection*{Step 11: Plug the concrete stepsize}
Choose $\eta_t=\frac{c}{\sqrt{t+1}}$. Then
\[
\sum_{t=0}^{T-1}\eta_t = c\sum_{t=0}^{T-1}\frac{1}{\sqrt{t+1}}
   \ge 2c\bigl(\sqrt{T}-1\bigr)=\Theta(\sqrt{T}),
\]
\[
\sum_{t=0}^{T-1}\eta_t^{2}=c^{2}\sum_{t=0}^{T-1}\frac{1}{t+1}
   \le c^{2}(1+\log T)=O(\log T).
\]
Insert into (Step 10):
\[
\E\!\bigl[f(\bar w_T)-f(w^\star)\bigr]
\le \frac{\norm{w_0-w^\star}^2}{2c\,\Theta(\sqrt{T})}
    +\frac{(1+q)G^{2}c^{2}O(\log T)}{2c\,\Theta(\sqrt{T})}
 = O\!\left(\frac{1}{\sqrt{T}}\right). \tag{Step 11}
\]

Thus Theorem~\ref{thm:conv} holds. \qed

\section*{Rate Derivation (With Constants)}
From (Step 10) we can write the bound explicitly:
\[
\E\!\bigl[f(\bar w_T)-f(w^\star)\bigr]
\le \frac{\norm{w_0-w^\star}^2}{2\sum_{t=0}^{T-1}\eta_t}
    +\frac{(1+q)G^{2}}{2}\,\frac{\sum_{t=0}^{T-1}\eta_t^{2}}{\sum_{t=0}^{T-1}\eta_t}.
\]
For $\eta_t=c/(t+1)^{\alpha}$ with $0.5\le\alpha<1$,
\[
\sum_{t=0}^{T-1}\eta_t = \Theta\bigl(T^{1-\alpha}\bigr),\qquad
\sum_{t=0}^{T-1}\eta_t^{2}= \Theta\bigl(T^{1-2\alpha}\bigr).
\]
Hence the rate becomes $O\!\bigl(T^{\alpha-1}\bigr)$; the optimal choice $\alpha=1/2$ yields the $O(1/\sqrt{T})$ rate.

\section*{Special Cases}
\begin{itemize}
\item \textbf{No quantization ($q=0$).} The bound reduces exactly to the classic SGD bound.
\item \textbf{Deterministic gradients ($\sigma=0$).} If $g_t=\nabla f(w_t)$, the term $G^{2}$ can be replaced by $\norm{\nabla f(w_t)}^{2}$, leading to tighter, data‑dependent rates.
\item \textbf{Strong convexity.} If $f$ is $\mu$‑strongly convex, a similar derivation with $\eta_t=\frac{2}{\mu (t+1)}$ gives
\[
\E\!\bigl[\norm{w_t-w^\star}^{2}\bigr]\le \frac{(1+q)G^{2}}{\mu^{2}t},
\]
i.e. an $O(1/t)$ rate.
\end{itemize}

\end{document}